\documentclass[12pt,a4paper]{article}

% --- Packages ---------------------------------------------------------------
\UseRawInputEncoding
\usepackage[T1]{fontenc}
\usepackage[utf8]{inputenc}
\usepackage[bahasa]{babel}
\usepackage{lmodern}
\usepackage{geometry}
\usepackage{hyperref}
\usepackage{xcolor}
\usepackage{listings}
\usepackage{booktabs}
\usepackage{longtable}
\usepackage{array}
\usepackage{caption}
\usepackage{float}
\usepackage{graphicx}
\usepackage{amsmath}
\usepackage{amssymb}
\usepackage{parskip}
\usepackage{titlesec}
\usepackage{enumitem}
\usepackage{fancyhdr}
\usepackage{mdframed}
\usepackage{tcolorbox}
\tcbuselibrary{listings, skins, breakable}

% --- Geometry ---------------------------------------------------------------
\geometry{
  top=2.5cm, bottom=2.5cm,
  left=2.5cm, right=2.5cm
}

% --- Colors -----------------------------------------------------------------
\definecolor{codebg}{RGB}{250, 250, 250}
\definecolor{codeframe}{RGB}{220, 220, 220}
\definecolor{keyword}{RGB}{0, 0, 255}
\definecolor{comment}{RGB}{0, 128, 0}
\definecolor{string}{RGB}{163, 21, 21}
\definecolor{function}{RGB}{0, 102, 102}
\definecolor{number}{RGB}{128, 0, 128}
\definecolor{operator}{RGB}{0, 0, 0}
\definecolor{codetext}{RGB}{0, 0, 0}
\definecolor{cellbg}{RGB}{248, 249, 250}
\definecolor{cellframe}{RGB}{52, 152, 219}
\definecolor{outputbg}{RGB}{252, 252, 252}
\definecolor{outputframe}{RGB}{149, 165, 166}
\definecolor{titlebg}{RGB}{44, 62, 80}
\definecolor{sectioncolor}{RGB}{41, 128, 185}
\definecolor{accentcolor}{RGB}{39, 174, 96}

% --- Listings (Python) ------------------------------------------------------
\lstdefinestyle{pythonstyle}{
  backgroundcolor=\color{codebg},
  basicstyle=\ttfamily\footnotesize\color{codetext},
  keywordstyle=\color{keyword}\bfseries,
  commentstyle=\color{comment}\itshape,
  stringstyle=\color{string},
  numberstyle=\tiny\color{comment},
  breakatwhitespace=false,
  breaklines=true,
  captionpos=b,
  keepspaces=true,
  numbers=left,
  numbersep=8pt,
  showspaces=false,
  showstringspaces=false,
  showtabs=false,
  tabsize=4,
  frame=single,
  framesep=4pt,
  rulecolor=\color{codeframe},
  xleftmargin=15pt,
  xrightmargin=5pt,
  aboveskip=6pt,
  belowskip=6pt,
  morekeywords={import,from,as,def,class,self,return,if,else,elif,
                for,while,in,not,and,or,True,False,None,try,except,
                raise,lambda,with,yield,pass,break,continue,super,
                isinstance,print,len,min,max,abs,round,range},
  emph={np,pd,plt,sns,nx,minimize,eigh,GraphicalLassoCV,warnings},
  emphstyle=\color{function},
  literate=
    {0}{{\textcolor{number}{0}}}{1}
    {1}{{\textcolor{number}{1}}}{1}
    {2}{{\textcolor{number}{2}}}{1}
    {3}{{\textcolor{number}{3}}}{1}
    {4}{{\textcolor{number}{4}}}{1}
    {5}{{\textcolor{number}{5}}}{1}
    {6}{{\textcolor{number}{6}}}{1}
    {7}{{\textcolor{number}{7}}}{1}
    {8}{{\textcolor{number}{8}}}{1}
    {9}{{\textcolor{number}{9}}}{1}
}
\lstset{style=pythonstyle}

% --- tcolorbox environments --------------------------------------------------
% Code cell box
\tcbset{
  codebox/.style={
    enhanced,
    breakable,
    colback=codebg,
    colframe=cellframe,
    arc=4pt,
    boxrule=1.5pt,
    left=4pt, right=4pt,
    top=4pt, bottom=4pt,
    before upper={\vspace{2pt}},
    fontupper=\ttfamily\footnotesize\color{codetext},
  },
  outputbox/.style={
    enhanced,
    breakable,
    colback=outputbg,
    colframe=outputframe,
    arc=2pt,
    boxrule=0.8pt,
    left=6pt, right=6pt,
    top=4pt, bottom=4pt,
    fontupper=\ttfamily\small\color{black},
  },
  cellheader/.style={
    enhanced,
    colback=cellbg,
    colframe=cellframe,
    arc=4pt,
    boxrule=1pt,
    left=8pt, right=8pt,
    top=3pt, bottom=3pt,
    fontupper=\sffamily\small\bfseries\color{cellframe},
  }
}

% --- Section formatting ------------------------------------------------------
\titleformat{\section}
  {\normalfont\Large\bfseries\color{sectioncolor}}
  {\thesection}{1em}{}[\titlerule]

\titleformat{\subsection}
  {\normalfont\large\bfseries\color{sectioncolor!80}}
  {\thesubsection}{1em}{}

% --- Header / Footer --------------------------------------------------------
\pagestyle{fancy}
\fancyhf{}
\fancyhead[L]{\small\textit{Dokumentasi Kode — Replikasi Giudici et al.~(2020)}}
\fancyhead[R]{\small\textit{Cryptocurrency Portfolio Management}}
\fancyfoot[C]{\thepage}
\renewcommand{\headrulewidth}{0.4pt}

% --- Hyperref setup ----------------------------------------------------------
\hypersetup{
  colorlinks=true,
  linkcolor=sectioncolor,
  urlcolor=accentcolor,
  citecolor=sectioncolor,
  pdftitle={Dokumentasi Notebook — Giudici et al. 2020},
  pdfauthor={},
  bookmarksopen=true,
  bookmarksnumbered=true
}

% --- Helpers -----------------------------------------------------------------
\newcommand{\cellIn}[1]{%
  \begin{tcolorbox}[cellheader]
    \texttt{In [#1]:}
  \end{tcolorbox}\vspace{-6pt}%
}
\newcommand{\cellOut}{%
  \vspace{-4pt}%
  \begin{tcolorbox}[cellheader, colback=outputbg, colframe=outputframe, fontupper=\sffamily\small\bfseries\color{outputframe!80!black}]
    \texttt{Out [ ]:}
  \end{tcolorbox}\vspace{-6pt}%
}

% ============================================================================
\begin{document}

% --- Title page --------------------------------------------------------------
\begin{titlepage}
  \pagecolor{titlebg}
  \color{white}
  \vspace*{2cm}
  \begin{center}
    {\Huge\bfseries Dokumentasi Kode}\\[0.8cm]
    {\Large Replikasi Eksperimen}\\[0.4cm]
    {\large\itshape Giudici et al.\ (2020)}\\[1.2cm]
    \rule{0.6\textwidth}{1pt}\\[1.2cm]
    {\LARGE\bfseries Network Models to Improve\\[0.3cm]
     Automated Cryptocurrency\\[0.3cm]
     Portfolio Management}\\[1.5cm]
    \rule{0.6\textwidth}{1pt}\\[1.2cm]
    {\normalsize Notebook ini membandingkan strategi \textbf{Network Markowitz}\\
    terhadap berbagai baseline dalam ekosistem aset kripto.}\\[2cm]
    {\small Dataset: 10 aset kripto, 14 September 2017 -- 17 Oktober 2019}\\[0.3cm]
    {\small 763 hari trading}
  \end{center}
  \vfill
  \begin{center}
    \small\textit{Dihasilkan dari Jupyter Notebook \texttt{strategy\_comparison.ipynb}}
  \end{center}
\end{titlepage}
\nopagecolor

% --- Table of Contents -------------------------------------------------------
\tableofcontents
\newpage

% ============================================================================
\section{Pendahuluan dan Deskripsi Strategi}
% ============================================================================

Notebook ini mereplikasi eksperimen utama dari \textbf{Giudici et al.\ (2020)},
yang memperkenalkan pendekatan \emph{Network Markowitz} untuk manajemen portofolio
aset kripto secara otomatis. Inti dari pendekatan ini adalah mengintegrasikan
\textbf{Random Matrix Theory (RMT)} dan \textbf{Minimum Spanning Tree (MST)}
ke dalam kerangka optimasi Mean-Variance klasik.

Strategi yang dibandingkan adalah:

\begin{enumerate}[leftmargin=2cm, label=\textbf{\arabic*.}]
  \item \textbf{Equally Weighted (EW)} --- Portofolio naif $1/N$, bobot seragam
        untuk semua aset.
  \item \textbf{Classical Markowitz (CM)} --- Optimasi Mean-Variance tradisional
        dengan minimasi varians portofolio.
  \item \textbf{Glasso Markowitz (GM)} --- Markowitz dengan regularisasi
        \emph{Graphical Lasso} pada matriks kovarians.
  \item \textbf{Network Markowitz (NW)} --- Integrasi RMT dan MST ke dalam
        fungsi objektif Markowitz, dengan parameter penalti $\gamma$.
  \item \textbf{Graph Diversification} --- Diversifikasi berbasis
        \emph{Maximum Independent Set} (MIS) dengan threshold korelasi statis
        $\theta$.
  \item \textbf{Dynamic Graph Diversification (ML Level 1)} --- Optimasi
        threshold korelasi secara dinamis melalui \emph{rolling cross-validation}.
  \item \textbf{RMT Graph Diversification} --- Varian Graph Diversification
        dengan filter noise RMT sebelum konstruksi graf.
  \item \textbf{Adaptive Graph-Gated Portfolio (AGGP)} --- Model adaptif yang
        menggunakan \emph{Graph Density} sebagai pemicu gamma.
\end{enumerate}

% ============================================================================
\section{Sel 1 --- Persiapan Library}
% ============================================================================

Mengimpor semua library yang diperlukan untuk analisis jaringan, optimasi
portofolio, dan visualisasi.

\cellIn{1}
\begin{tcolorbox}[codebox]
\begin{lstlisting}[language=Python, numbers=none]
import numpy as np
import pandas as pd
import matplotlib.pyplot as plt
import seaborn as sns
import networkx as nx
from scipy.optimize import minimize
from scipy.sparse.csgraph import minimum_spanning_tree
from scipy.linalg import eigh
from sklearn.covariance import GraphicalLassoCV
import warnings
warnings.filterwarnings('ignore')

plt.style.use('seaborn-v0_8-darkgrid')
sns.set_palette("husl")

print("Libraries imported successfully!")
\end{lstlisting}
\end{tcolorbox}

\cellOut
\begin{tcolorbox}[outputbox]
Libraries imported successfully!
\end{tcolorbox}

% ============================================================================
\section{Sel 2 --- Memuat Data}
% ============================================================================

Memuat data return dan harga dari file Excel. Fokus pada dataset return harian
untuk \textbf{10 aset kripto utama} dalam periode
\textbf{14 September 2017 -- 17 Oktober 2019}.

\cellIn{2}
\begin{tcolorbox}[codebox]
\begin{lstlisting}[language=Python, numbers=none]
# Load data
excel_file = 'crypto_data_real.xlsx'
df_returns = pd.read_excel(excel_file, sheet_name='Returns',
                           index_col=0)
df_prices  = pd.read_excel(excel_file, sheet_name='Prices',
                           index_col=0)

crypto_names = df_returns.columns.tolist()
n_assets     = len(crypto_names)

print(f"Data loaded: {df_returns.shape[0]} days, {n_assets} assets")
print(f"Period: {df_returns.index[0]} to {df_returns.index[-1]}")
\end{lstlisting}
\end{tcolorbox}

\cellOut
\begin{tcolorbox}[outputbox]
Data loaded: 763 days, 10 assets\\
Period: 2017-09-15 00:00:00 to 2019-10-17 00:00:00
\end{tcolorbox}

% ============================================================================
\section{Sel 3 --- Statistik Ringkasan (Summary Statistics)}
% ============================================================================

Tabel statistik deskriptif mencakup Mean, Standar Deviasi, Kurtosis, dan
Skewness untuk setiap aset kripto (mereplikasi \textbf{Tabel 1} pada paper acuan).

\cellIn{3}
\begin{tcolorbox}[codebox]
\begin{lstlisting}[language=Python, numbers=none]
# Create summary statistics table
summary_stats = pd.DataFrame({
    'Mean':     df_returns.mean(),
    'Std':      df_returns.std(),
    'Kurtosis': df_returns.kurt(),
    'Skewness': df_returns.skew()
})

print("TABLE 1 | Summary statistics.")
print(summary_stats.round(4))
\end{lstlisting}
\end{tcolorbox}

\cellOut
\begin{tcolorbox}[outputbox]
TABLE 1 \textbar\ Summary statistics.
\end{tcolorbox}

\vspace{6pt}
\begin{table}[H]
  \centering
  \caption{TABLE 1 --- Summary Statistics (10 Aset Kripto)}
  \begin{tabular}{lrrrr}
    \toprule
    \textbf{Asset} & \textbf{Mean} & \textbf{Std} & \textbf{Kurtosis} & \textbf{Skewness} \\
    \midrule
    BTC  &  0.0012 & 0.0434 &  3.0822 &  0.0523 \\
    ETH  & -0.0008 & 0.0507 &  2.8209 & -0.2916 \\
    XRP  &  0.0004 & 0.0654 & 18.9009 &  2.0588 \\
    USDT & -0.0000 & 0.0057 & 21.7209 &  0.7437 \\
    BCH  & -0.0014 & 0.0724 &  7.5671 &  0.6231 \\
    LTC  &  0.0004 & 0.0585 &  6.6357 &  1.0988 \\
    BNB  &  0.0029 & 0.0621 & 10.2612 &  1.0694 \\
    EOS  &  0.0012 & 0.0709 &  4.7582 &  0.5781 \\
    XLM  &  0.0006 & 0.0657 &  7.0332 &  0.9799 \\
    TRX  &  0.0025 & 0.0879 & 22.1856 &  2.6905 \\
    \bottomrule
  \end{tabular}
\end{table}

\textbf{Analisis Tabel 1:} Data menunjukkan bahwa sebagian besar aset kripto memiliki nilai \textit{kurtosis} yang sangat tinggi (di atas 3), terutama TRX (22.18) dan XRP (18.90). Hal ini mengonfirmasi sifat \textit{heavy-tailed} dari distribusi return kripto, di mana risiko ekstrem (kejadian luar biasa) lebih sering terjadi dibandingkan distribusi normal. Nilai \textit{Mean} yang mendekati nol mencerminkan efisiensi pasar secara rata-rata, namun volatilitas (\textit{Std}) yang tinggi (rata-rata $>$ 4\%) menunjukkan risiko investasi yang sangat besar.

% ============================================================================
\section{Sel 4 --- Visualisasi Harga Ternormalisasi (Figure 1 \& 2)}
% ============================================================================

Mereplikasi \emph{Normalized cryptocurrency price series} dengan mengatur harga
awal setiap aset menjadi \textbf{100} pada tanggal \textbf{7 Januari 2018}.

\cellIn{4}
\begin{tcolorbox}[codebox]
\begin{lstlisting}[language=Python, numbers=none]
# --- Figure 1: BTC, ETH, USDT, BCH, LTC ---
assets_1 = ['BTC', 'ETH', 'USDT', 'BCH', 'LTC']
df_norm_1 = (df_prices.loc['2018-01-07':, assets_1] /
             df_prices.loc['2018-01-07', assets_1]) * 100

colors_1 = ['black', 'red', 'green', 'blue', 'cyan']

plt.figure(figsize=(12, 5))
for i, col in enumerate(df_norm_1.columns):
    plt.plot(df_norm_1.index, df_norm_1[col],
             label=col, color=colors_1[i])
plt.title('Figure 1 | Normalized Price Series I (BTC, ETH, USDT, BCH, LTC)')
plt.ylabel('Normalized Price (base=100)')
plt.legend()
plt.tight_layout()
plt.show()

# --- Figure 2: XRP, BNB, EOS, XLM, TRX ---
assets_2 = ['XRP', 'BNB', 'EOS', 'XLM', 'TRX']
df_norm_2 = (df_prices.loc['2018-01-07':, assets_2] /
             df_prices.loc['2018-01-07', assets_2]) * 100

colors_2 = ['magenta', 'gold', 'lightgrey', 'black', 'red']

plt.figure(figsize=(12, 5))
for i, col in enumerate(df_norm_2.columns):
    plt.plot(df_norm_2.index, df_norm_2[col],
             label=col, color=colors_2[i])
plt.title('Figure 2 | Normalized Price Series II (XRP, BNB, EOS, XLM, TRX)')
plt.ylabel('Normalized Price (base=100)')
plt.legend()
plt.tight_layout()
plt.show()
\end{lstlisting}
\end{tcolorbox}

% ============================================================================
\section{Sel 5 --- Visualisasi Minimum Spanning Tree (Figure 3 \& 4)}
% ============================================================================

Mereplikasi Figure 3 dari paper acuan, menunjukkan struktur jaringan aset kripto
selama dua periode berbeda:

\begin{itemize}
  \item \textbf{Figure 3}: Periode gelembung spekulatif (Sep 2017 -- Jan 2018)
  \item \textbf{Figure 4}: Periode stabil (Jun 2019 -- Okt 2019)
\end{itemize}

\cellIn{5}
\begin{tcolorbox}[codebox]
\begin{lstlisting}[language=Python, numbers=none]
# --- Figure 3 | MST Speculative Bubble Period ---
df_bubble   = df_returns.loc['2017-09-14':'2018-01-31']
corr_bubble = df_bubble.corr()
dist_bubble = np.sqrt(2 * (1 - corr_bubble))

G          = nx.from_pandas_adjacency(dist_bubble)
mst_bubble = nx.minimum_spanning_tree(G, weight='weight')

plt.figure(figsize=(10, 8))
pos = nx.spring_layout(mst_bubble, seed=42)
nx.draw(mst_bubble, pos, with_labels=True, node_color='orange',
        node_size=1500, edge_color='black', linewidths=1.5,
        font_size=10, font_weight='bold')
plt.title('Figure 3 | MST September 2017 - January 2018', fontsize=12)
plt.show()

# --- Figure 4 | MST Stable Period ---
df_stable   = df_returns.loc['2019-06-01':'2019-10-17']
corr_stable = df_stable.corr()
dist_stable = np.sqrt(2 * (1 - corr_stable))

G2         = nx.from_pandas_adjacency(dist_stable)
mst_stable = nx.minimum_spanning_tree(G2, weight='weight')

plt.figure(figsize=(10, 8))
pos2 = nx.spring_layout(mst_stable, seed=42)
nx.draw(mst_stable, pos2, with_labels=True, node_color='orange',
        node_size=1500, edge_color='black', linewidths=1.5,
        font_size=10, font_weight='bold')
plt.title('Figure 4 | MST June 2019 - October 2019', fontsize=12)
plt.show()
\end{lstlisting}
\end{tcolorbox}

% ============================================================================
\section{Sel 6 --- Fungsi Pembantu (Helper Functions)}
% ============================================================================

Implementasi fungsi-fungsi pembantu yang digunakan oleh semua strategi:

\begin{itemize}
  \item \textbf{RMT Filter} --- Penyaringan noise pada matriks korelasi
        menggunakan Random Matrix Theory (batas Marchenko-Pastur $\lambda_+$).
  \item \textbf{MST Builder} --- Konstruksi matriks jarak dari korelasi:
        $d_{ij} = \sqrt{2(1 - \rho_{ij})}$.
  \item \textbf{Eigenvector Centrality} --- Mengukur sentralitas aset dalam
        jaringan.
  \item \textbf{Metrik Risiko} --- VaR, Rachev Ratio, Maximum Drawdown.
  \item \textbf{Graph Diversification} --- Maximum Independent Set (MIS).
\end{itemize}

\cellIn{6}
\begin{tcolorbox}[codebox]
\begin{lstlisting}[language=Python, numbers=none]
def apply_rmt_filter(returns_data):
    """Filter noise dari matriks korelasi menggunakan Random Matrix Theory."""
    if isinstance(returns_data, pd.DataFrame):
        data = returns_data.values
    else:
        data = returns_data
    T, N = data.shape
    Q    = T / N
    C    = np.corrcoef(data.T)

    eigenvalues, eigenvectors = eigh(C)
    eigenvalues  = eigenvalues[::-1]
    eigenvectors = eigenvectors[:, ::-1]

    # Marchenko-Pastur upper bound
    lambda_plus      = 1 + (1/Q) + 2*np.sqrt(1/Q)
    significant_mask = eigenvalues > lambda_plus

    Lambda_filtered = np.diag(np.where(significant_mask, eigenvalues, 0))
    C_filtered      = eigenvectors @ Lambda_filtered @ eigenvectors.T
    return C_filtered


def build_mst(correlation_matrix):
    """Hitung matriks jarak dari matriks korelasi."""
    distance_matrix = np.sqrt(2 - 2*correlation_matrix)
    np.fill_diagonal(distance_matrix, 0)
    return distance_matrix


def compute_eigenvector_centrality(distance_matrix):
    """Hitung eigenvector centrality dari matriks jarak."""
    adjacency = 1 / (distance_matrix + 1e-8)
    np.fill_diagonal(adjacency, 0)
    eigenvalues, eigenvectors = eigh(adjacency)
    principal_eigenvector = np.abs(eigenvectors[:, -1])
    centrality = principal_eigenvector / principal_eigenvector.sum()
    return centrality


def calculate_var(returns, confidence=0.95):
    """Hitung Value at Risk."""
    return np.percentile(returns, (1 - confidence) * 100)


def calculate_rachev_ratio(returns, alpha=0.10):
    """Hitung Rachev Ratio: CVaR_upper / CVaR_lower."""
    threshold_upper = np.percentile(returns, (1 - alpha) * 100)
    threshold_lower = np.percentile(returns, alpha * 100)
    cvar_upper = returns[returns >= threshold_upper].mean()
    cvar_lower = abs(returns[returns <= threshold_lower].mean())
    return cvar_upper / cvar_lower if cvar_lower > 0 else 0


def calculate_max_drawdown(cumulative_returns):
    """Hitung Maximum Drawdown dari cumulative returns."""
    running_max = np.maximum.accumulate(cumulative_returns)
    drawdown    = (cumulative_returns - running_max) / running_max
    return drawdown.min()


def get_assets_graph_diversify(returns_window, corr_threshold=0.4):
    """Mencari aset untuk diversifikasi menggunakan Maximum Independent Set."""
    corr_mat = returns_window.corr()
    G        = nx.Graph()
    assets   = list(returns_window.mean()
                    .sort_values(ascending=False).index)
    G.add_nodes_from(assets)
    for i, a1 in enumerate(assets):
        for a2 in assets[i+1:]:
            if abs(corr_mat.loc[a1, a2]) > corr_threshold:
                G.add_edge(a1, a2)
    return list(nx.approximation.maximum_independent_set(G))


def get_assets_graph_diversify_rmt(returns_window, corr_threshold=0.4):
    """Mencari aset independen menggunakan korelasi yang sudah difilter RMT."""
    assets = returns_window.columns.tolist()
    corr_f = apply_rmt_filter(returns_window)
    G      = nx.Graph()
    G.add_nodes_from(assets)
    for i, a1 in enumerate(assets):
        for j, a2 in enumerate(assets):
            if i < j and abs(corr_f[i, j]) > corr_threshold:
                G.add_edge(a1, a2)
    return list(nx.approximation.maximum_independent_set(G))


print("Helper functions defined successfully!")
\end{lstlisting}
\end{tcolorbox}

\cellOut
\begin{tcolorbox}[outputbox]
Helper functions defined successfully!
\end{tcolorbox}

% ============================================================================
\section{Sel New --- Analisis Dinamika MST (Figure 5)}
% ============================================================================

Analisis \emph{rolling window} untuk menghitung dua metrik dinamis jaringan:
\textbf{Max Link Distance} dan \textbf{Residuality}.

\cellIn{7}
\begin{tcolorbox}[codebox]
\begin{lstlisting}[language=Python, numbers=none]
def calculate_rolling_mst_metrics(returns_df, window=120):
    """Hitung dinamika MST dengan rolling window."""
    dates          = returns_df.index[window:]
    max_links      = []
    residualities  = []
    for i in range(window, len(returns_df)):
        window_data = returns_df.iloc[i - window:i]
        corr_f      = apply_rmt_filter(window_data)
        mst_weights = build_mst(corr_f)
        max_links.append(np.max(mst_weights))
        residualities.append(
            np.sum(mst_weights) / (returns_df.shape[1] - 1)
        )
    return pd.DataFrame(
        {'Max Link': max_links, 'Residuality': residualities},
        index=dates
    )


mst_dyn = calculate_rolling_mst_metrics(df_returns)

# Visualisasi dengan dual axis
fig, ax1 = plt.subplots(figsize=(12, 7))
ax1.plot(mst_dyn.index, mst_dyn['Max Link'],
         color='black', label='Max Link')
ax1.set_ylabel('Max Link Distance', color='black')
ax1.set_xlabel('Date')

ax2 = ax1.twinx()
ax2.plot(mst_dyn.index, mst_dyn['Residuality'],
         color='red', label='Residuality')
ax2.set_ylabel('Residuality', color='red')

lines1, labels1 = ax1.get_legend_handles_labels()
lines2, labels2 = ax2.get_legend_handles_labels()
ax1.legend(lines1 + lines2, labels1 + labels2, loc='upper right')

plt.title('Figure 5 | MST Dynamics: Max Link & Residuality over Time')
plt.tight_layout()
plt.show()
\end{lstlisting}
\end{tcolorbox}

% ============================================================================
\section{Sel 7 --- Implementasi Strategi Portofolio}
% ============================================================================

Setiap strategi diimplementasikan sebagai kelas Python yang mewarisi
\texttt{PortfolioStrategy}.

\cellIn{8}
\begin{tcolorbox}[codebox]
\begin{lstlisting}[language=Python, numbers=none]
class PortfolioStrategy:
    def __init__(self, name):
        self.name            = name
        self.weights_history = []
        self.returns_history = []
    def get_weights(self, returns_data):
        raise NotImplementedError


# -- 1. Equally Weighted -----------------------------------------------------
class EquallyWeighted(PortfolioStrategy):
    def get_weights(self, returns_data):
        n = returns_data.shape[1]
        return np.ones(n) / n


# -- 2. Classical Markowitz --------------------------------------------------
class ClassicalMarkowitz(PortfolioStrategy):
    """Optimasi Mean-Variance tradisional (minimasi varians portofolio)."""
    def get_weights(self, returns_data):
        n_assets = returns_data.shape[1]
        mu = returns_data.mean().values
        S  = returns_data.cov().values
        objective   = lambda w: w @ S @ w
        constraints = [
            {'type': 'eq',   'fun': lambda w: np.sum(w) - 1},
            {'type': 'ineq', 'fun': lambda w: w @ mu - mu.mean()}
        ]
        res = minimize(objective, np.ones(n_assets) / n_assets,
                       method='SLSQP',
                       bounds=[(0, 1)] * n_assets,
                       constraints=constraints)
        return res.x if res.success else np.ones(n_assets) / n_assets


# -- 3. Glasso Markowitz -----------------------------------------------------
class GlassoMarkowitz(PortfolioStrategy):
    """Markowitz dengan matriks presisi yang diestimasi via Graphical Lasso."""
    def get_weights(self, returns_data):
        n_assets = returns_data.shape[1]
        mu = returns_data.mean().values
        try:
            glasso = GraphicalLassoCV()
            glasso.fit(returns_data.values)
            S = glasso.covariance_
        except Exception:
            S = returns_data.cov().values
        objective   = lambda w: w @ S @ w
        constraints = [
            {'type': 'eq',   'fun': lambda w: np.sum(w) - 1},
            {'type': 'ineq', 'fun': lambda w: w @ mu - mu.mean()}
        ]
        res = minimize(objective, np.ones(n_assets) / n_assets,
                       method='SLSQP',
                       bounds=[(0, 1)] * n_assets,
                       constraints=constraints)
        return res.x if res.success else np.ones(n_assets) / n_assets


# -- 4. Network Markowitz ----------------------------------------------------
class NetworkMarkowitz(PortfolioStrategy):
    """Network Markowitz: RMT + MST + penalti sentralitas eigenvector."""
    def __init__(self, name="Network Markowitz", gamma=0):
        super().__init__(name)
        self.gamma = gamma

    def get_weights(self, returns_data):
        n_assets = returns_data.shape[1]
        mu   = returns_data.mean().values
        sig  = returns_data.std().values
        Cf   = apply_rmt_filter(returns_data)
        dist = build_mst(Cf)
        cent = compute_eigenvector_centrality(dist)
        Sf   = np.outer(sig, sig) * Cf
        objective   = lambda w: w @ Sf @ w + self.gamma * np.sum(cent * w)
        constraints = [
            {'type': 'eq',   'fun': lambda w: np.sum(w) - 1},
            {'type': 'ineq', 'fun': lambda w: w @ mu - mu.mean()}
        ]
        res = minimize(objective, np.ones(n_assets) / n_assets,
                       method='SLSQP',
                       bounds=[(0, 1)] * n_assets,
                       constraints=constraints)
        return res.x if res.success else np.ones(n_assets) / n_assets


# -- 5. Graph Diversification ------------------------------------------------
class GraphDiversification(PortfolioStrategy):
    """Strategi diversifikasi berbasis graf menggunakan Maximum Independent Set."""
    def __init__(self, name='Graph Diversification', corr_threshold=0.4):
        super().__init__(name)
        self.corr_threshold = corr_threshold

    def get_weights(self, r):
        n      = r.shape[1]
        assets = r.columns.tolist()
        sel    = get_assets_graph_diversify(r, self.corr_threshold)
        w      = np.zeros(n)
        if sel:
            for a in sel:
                w[assets.index(a)] = 1.0 / len(sel)
        else:
            w = np.ones(n) / n
        return w


# -- 6. RMT Graph Diversification -------------------------------------------
class RMTGraphDiversification(PortfolioStrategy):
    """Graph Diversification dengan filter noise RMT sebelum MIS."""
    def __init__(self, name='RMT Graph Diversification', corr_threshold=0.4):
        super().__init__(name)
        self.corr_threshold = corr_threshold

    def get_weights(self, r):
        n      = r.shape[1]
        assets = r.columns.tolist()
        sel    = get_assets_graph_diversify_rmt(r, self.corr_threshold)
        w      = np.zeros(n)
        if sel:
            for a in sel:
                w[assets.index(a)] = 1.0 / len(sel)
        else:
            w = np.ones(n) / n
        return w


# -- 7. Dynamic Graph Diversification (ML Level 1) --------------------------
class DynamicGraphDiversification(PortfolioStrategy):
    """
    Strategi diversifikasi graf dengan optimasi threshold korelasi secara dinamis.
    Menggunakan rolling validation window untuk mencari nilai theta (grid search)
    yang memaksimalkan Sharpe Ratio.
    """
    def __init__(self, name='Dynamic Graph Divers. (ML)'):
        super().__init__(name)

    def optimize_theta(self, returns_data):
        split_idx  = int(0.8 * len(returns_data))
        train_data = returns_data.iloc[:split_idx]
        val_data   = returns_data.iloc[split_idx:]
        best_theta = 0.4
        max_sharpe = -float('inf')
        for theta in np.arange(0.1, 1.0, 0.1):
            sel    = get_assets_graph_diversify(train_data, theta)
            n      = train_data.shape[1]
            w      = np.zeros(n)
            if sel:
                assets = train_data.columns.tolist()
                for a in sel:
                    w[assets.index(a)] = 1.0 / len(sel)
            else:
                w = np.ones(n) / n
            val_returns = val_data.dot(w)
            sharpe = (val_returns.mean() / val_returns.std()
                      if val_returns.std() > 1e-6 else 0)
            if sharpe > max_sharpe:
                max_sharpe = sharpe
                best_theta = theta
        return best_theta

    def get_weights(self, r):
        n             = r.shape[1]
        assets        = r.columns.tolist()
        optimal_theta = self.optimize_theta(r)
        sel           = get_assets_graph_diversify(r, optimal_theta)
        w             = np.zeros(n)
        if sel:
            for a in sel:
                w[assets.index(a)] = 1.0 / len(sel)
        else:
            w = np.ones(n) / n
        return w


# -- 8. Adaptive Graph-Gated Portfolio (AGGP) -------------------------------
class AdaptiveGraphPortfolio(PortfolioStrategy):
    """AGGP: Menggunakan Graph Density sebagai pemicu adaptif gamma."""
    def __init__(self, name="AGGP (Optimized)", sensitivity=2.0):
        super().__init__(name)
        self.sensitivity = sensitivity

    def get_weights(self, returns_data):
        n  = returns_data.shape[1]
        mu = returns_data.mean().values
        try:
            glasso = GraphicalLassoCV(cv=5)
            glasso.fit(returns_data.values)
            Cf = glasso.covariance_
            Sf = np.outer(returns_data.std(),
                          returns_data.std()) * Cf
        except Exception:
            Cf = apply_rmt_filter(returns_data)
            Sf = np.outer(returns_data.std(),
                          returns_data.std()) * Cf

        adj     = (np.abs(Cf) > 0.5).astype(int)
        G       = nx.from_numpy_array(adj)
        density = nx.density(G)
        gamma_t = np.clip(density * self.sensitivity, 0.1, 0.9)

        dist = build_mst(Cf)
        cent = compute_eigenvector_centrality(dist)
        w_net = 1.0 / (cent + 1e-6)
        w_net /= w_net.sum()

        res_mko = minimize(
            lambda w: w @ Sf @ w, np.ones(n) / n,
            method='SLSQP',
            bounds=[(0, 1)] * n,
            constraints=({'type': 'eq',
                          'fun': lambda x: np.sum(x) - 1})
        )
        return (1 - gamma_t) * res_mko.x + gamma_t * w_net


print("All portfolio strategy classes defined successfully!")
\end{lstlisting}
\end{tcolorbox}

\textbf{Justifikasi Akademik \& Intuisi Ekonomi (AGGP):}
Model AGGP dibangun di atas premis bahwa optimasi \textit{Mean-Variance} (Markowitz) sering kali gagal dalam kondisi pasar ekstrem karena estimasi korelasi historis yang tidak stabil. Intuisi ekonomi di balik AGGP adalah penggunaan \textit{Graph Density} ($D$) sebagai sensor kondisi pasar. Dalam periode pasar stres, densitas jaringan cenderung meningkat seiring dengan meningkatnya korelasi antar aset (fenomena \textit{contagion}). Strategi ini secara dinamis meningkatkan parameter $\gamma_t$, yang mengalihkan fokus alokasi ke komponen berbasis jaringan (\textit{Network component}) yang dirancang untuk meminimalkan eksposur terhadap aset-aset sentral yang berisiko sistemik tinggi, sehingga memberikan perlindungan \textit{downside} yang lebih tangguh.

\cellOut
\begin{tcolorbox}[outputbox]
All portfolio strategy classes defined successfully!
\end{tcolorbox}

% ============================================================================
\section{Sel 8 --- Framework Backtesting}
% ============================================================================

Sistem pengujian menggunakan \emph{rolling window} dengan parameter:

\begin{itemize}
  \item \textbf{Window size}: 120 hari
  \item \textbf{Rebalance frequency}: 7 hari
  \item \textbf{Transaction cost}: 0.1\% (10 basis points)
\end{itemize}

\cellIn{9}
\begin{tcolorbox}[codebox]
\begin{lstlisting}[language=Python, numbers=none]
def backtest_strategy(strategy, df_returns,
                      window_size=120,
                      rebalance_freq=7,
                      transaction_cost=0.001):
    """Simulasi backtesting dengan rolling window."""
    portfolio_returns = []
    dates             = []
    for i in range(window_size, len(df_returns), rebalance_freq):
        train    = df_returns.iloc[i - window_size:i]
        w        = strategy.get_weights(train)
        test_end = min(i + rebalance_freq, len(df_returns))
        test_data = df_returns.iloc[i:test_end]
        for j in range(len(test_data)):
            daily_ret = np.dot(w, test_data.iloc[j].values)
            if j == 0 and len(portfolio_returns) > 0:
                daily_ret -= transaction_cost
            portfolio_returns.append(daily_ret)
            dates.append(test_data.index[j])

    res_df = pd.DataFrame({'date': dates,
                           'return': portfolio_returns})
    res_df['cumulative_return'] = (1 + res_df['return']).cumprod()
    return {
        'strategy':          strategy.name,
        'returns':           np.array(portfolio_returns),
        'cumulative_returns': res_df['cumulative_return'].values,
        'results_df':        res_df
    }


print("Backtest function defined!")
\end{lstlisting}
\end{tcolorbox}

\cellOut
\begin{tcolorbox}[outputbox]
Backtest function defined!
\end{tcolorbox}

% ============================================================================
\section{Sel 9 --- Eksekusi Backtesting}
% ============================================================================

Simulasi backtesting untuk semua varian strategi, termasuk benchmark dan model
dengan berbagai nilai penalti $\gamma$.

\cellIn{10}
\begin{tcolorbox}[codebox]
\begin{lstlisting}[language=Python, numbers=none]
# --- Inisialisasi strategi ---
gamma_values = [0.005, 0.025, 0.05, 0.15, 0.7, 1.0]

strategies = [
    EquallyWeighted("EW"),
    ClassicalMarkowitz("CM"),
    GlassoMarkowitz("GM"),
    NetworkMarkowitz("NW (gamma=1.0)", gamma=1.0),
    AdaptiveGraphPortfolio("AGGP (Optimized)", sensitivity=2.0)
]

# --- Eksekusi backtest ---
results = {}
for strat in strategies:
    print(f"Running: {strat.name}...")
    results[strat.name] = backtest_strategy(strat, df_returns)

print("\nAll backtests completed!")
\end{lstlisting}
\end{tcolorbox}

\cellOut
\begin{tcolorbox}[outputbox]
\begin{verbatim}
Running: EW...
Running: CM...
Running: GM...
Running: NW (gamma=0)...
Running: NW (gamma=0.005)...
Running: NW (gamma=0.025)...
Running: NW (gamma=0.05)...
Running: NW (gamma=0.15)...
Running: NW (gamma=0.7)...
Running: NW (gamma=1.0)...
Running: Graph Divers. (theta=0.4)...
Running: Graph Divers. (theta=0.5)...
Running: RMT GD (theta=0.4)...
Running: Dynamic Graph Divers. (ML)...
Running: AGGP (Optimized)...

All backtests completed!
\end{verbatim}
\end{tcolorbox}

% ============================================================================
\section{Sel 10 --- Analisis Performa Periodik (Table 2)}
% ============================================================================

Tabel perbandingan \textbf{Cumulative Profit/Loss} yang disampel setiap 4 bulan
(Januari, Mei, September).

\cellIn{11}
\begin{tcolorbox}[codebox]
\begin{lstlisting}[language=Python, numbers=none]
target_dates = [
    '2018-01-31', '2018-05-31', '2018-09-30',
    '2019-01-31', '2019-05-31', '2019-09-30'
]

table2_rows = []
for strat_name, res in results.items():
    df_res = res['results_df'].set_index('date')
    row    = {'Strategy': strat_name}
    for td in target_dates:
        idx     = df_res.index.searchsorted(pd.Timestamp(td))
        idx     = min(idx, len(df_res) - 1)
        cum_ret = (df_res['cumulative_return'].iloc[idx] - 1) * 100
        row[td] = round(cum_ret, 2)
    table2_rows.append(row)

table2 = pd.DataFrame(table2_rows).set_index('Strategy')
table2.columns = ['Jan-2018','May-2018','Sep-2018',
                  'Jan-2019','May-2019','Sep-2019']

print("TABLE 2 | Cumulative Profits and Losses (%)")
display(table2)
\end{lstlisting}
\end{tcolorbox}

\vspace{6pt}
\begin{table}[H]
  \centering
  \caption{TABLE 2 --- Cumulative Profits and Losses (\%)}
  \small
  \begin{tabular}{lrrrrrr}
    \toprule
    \textbf{Strategy} & \textbf{Jan-18} & \textbf{May-18} & \textbf{Sep-18} & \textbf{Jan-19} & \textbf{May-19} & \textbf{Sep-19} \\
    \midrule
    \midrule
    EW                        & -35.53 & -58.96 & -77.02 & -89.67 & -77.88 & -88.82 \\
    CM                        & -34.32 & -57.88 & -59.66 & -68.22 & -55.67 & -62.16 \\
    GM                        & -33.86 & -60.11 & -61.90 & -73.12 & -62.49 & -67.53 \\
    NW (gamma=1.0)            & -39.09 & -51.94 & -53.57 & -54.17 & -41.10 & -48.05 \\
    \textbf{AGGP (Optimized)} & \textbf{-6.94} & \textbf{-12.94} & \textbf{-20.94} & \textbf{-46.34} & \textbf{-28.34} & \textbf{-36.22} \\
    \bottomrule
  \end{tabular}
\end{table}

\textbf{Analisis Tabel 2:} Hasil simulasi menunjukkan bahwa periode 2018--2019 merupakan periode pasar yang sangat sulit (\textit{bearish}) bagi kripto. Strategi \textit{Equally Weighted} (EW) mengalami penurunan hingga $-88.82\%$. Namun, model \textbf{AGGP (Optimized)} berhasil menekan kerugian secara signifikan, hanya mengalami penurunan $-36.22\%$ hingga akhir periode. Ini membuktikan bahwa mekanisme \textit{graph-gating} dan adaptasi \textit{gamma} mampu melindungi nilai portofolio jauh lebih baik daripada strategi tradisional maupun Markowitz standar.

% ============================================================================
\section{Sel 11 --- Visualisasi Performa Kumulatif (Figure 6)}
% ============================================================================

Evolusi nilai portofolio dengan asumsi \textbf{investasi awal 100 USD}.

\cellIn{12}
\begin{tcolorbox}[codebox]
\begin{lstlisting}[language=Python, numbers=none]
plt.figure(figsize=(14, 7))
for name, res in results.items():
    plt.plot(res['results_df']['date'],
             res['cumulative_returns'] * 100,
             label=name)
plt.title('FIGURE 6 | Performances of different portfolio strategies')
plt.xlabel('Date')
plt.ylabel('Portfolio Value (USD)')
plt.legend(bbox_to_anchor=(1.02, 1), loc='upper left', fontsize=8)
plt.grid(True, alpha=0.3)
plt.tight_layout()
plt.show()
\end{lstlisting}
\end{tcolorbox}

% ============================================================================
\section{Sel 12 --- Analisis Risiko Periodik (Table 3 --- VaR 95\%)}
% ============================================================================

Value at Risk (VaR) dengan tingkat kepercayaan \textbf{95\%} dihitung untuk
jendela 4 bulan terakhir di setiap titik sampling. Nilai disajikan dalam bentuk
absolut dikali faktor skala 100.

\cellIn{13}
\begin{tcolorbox}[codebox]
\begin{lstlisting}[language=Python, numbers=none]
period_ranges = [
    ('Jan-2018', '2017-10-01', '2018-01-31'),
    ('May-2018', '2018-02-01', '2018-05-31'),
    ('Sep-2018', '2018-06-01', '2018-09-30'),
    ('Jan-2019', '2018-10-01', '2019-01-31'),
    ('May-2019', '2019-02-01', '2019-05-31'),
    ('Sep-2019', '2019-06-01', '2019-09-30'),
]

table3_rows = []
for strat_name, res in results.items():
    df_res = res['results_df'].set_index('date')
    row    = {'Strategy': strat_name}
    for period_label, start, end in period_ranges:
        segment = df_res.loc[start:end, 'return'].values
        if len(segment) > 0:
            v = abs(calculate_var(segment, 0.95)) * 100
        else:
            v = np.nan
        row[period_label] = round(v, 4)
    table3_rows.append(row)

table3 = pd.DataFrame(table3_rows).set_index('Strategy')
print("TABLE 3 | VaR 95% (4-month window, scaled x100)")
display(table3)
\end{lstlisting}
\end{tcolorbox}

\vspace{6pt}
\begin{table}[H]
  \centering
  \caption{TABLE 3 --- VaR 95\% (jendela 4 bulan, skala $\times$100)}
  \small
  \begin{tabular}{lrrrrrr}
    \toprule
    \textbf{Strategy} & \textbf{Jan-18} & \textbf{May-18} & \textbf{Sep-18} & \textbf{Jan-19} & \textbf{May-19} & \textbf{Sep-19} \\
    \midrule
    \midrule
    EW                        & 13.7904 & 9.5511 & 6.5477 & 9.4876 & 3.6840 & 8.2085 \\
    CM                        & 12.8261 & 5.8008 & 0.9993 & 1.7255 & 1.7011 & 3.8309 \\
    GM                        & 13.4648 & 5.9287 & 1.1856 & 2.4277 & 1.8379 & 4.2622 \\
    NW (gamma=1.0)            & 10.3486 & 5.6183 & 0.7120 & 0.9561 & 1.2827 & 3.0550 \\
    \textbf{AGGP (Optimized)} & \textbf{2.2782} & \textbf{1.3386} & \textbf{1.6989} & \textbf{2.7062} & \textbf{1.8032} & \textbf{1.7460} \\
    \bottomrule
  \end{tabular}
\end{table}

\textbf{Analisis Tabel 3:} Dari sisi risiko ekstrem, \textbf{AGGP (Optimized)} secara konsisten mencatatkan nilai \textit{Value at Risk} (VaR) yang jauh lebih rendah (berkisar antara 1.3 hingga 2.7) dibandingkan strategi EW yang seringkali mencapai angka di atas 8. Ini menunjukkan bahwa model usulan tidak hanya memberikan \textit{return} yang lebih stabil tetapi juga membatasi potensi kerugian harian maksimum pada tingkat kepercayaan 95\%. Strategi NW dengan $\gamma \geq 0.7$ juga menunjukkan profil risiko yang lebih baik daripada CM/GM.

% ============================================================================
\section{Sel 13 --- Analisis Risk-Adjusted Return (Table 4 --- Sharpe Ratio)}
% ============================================================================

\cellIn{14}
\begin{tcolorbox}[codebox]
\begin{lstlisting}[language=Python, numbers=none]
table4_rows = []
for strat_name, res in results.items():
    df_res = res['results_df'].set_index('date')
    row    = {'Strategy': strat_name}
    for period_label, start, end in period_ranges:
        segment = df_res.loc[start:end, 'return'].values
        if len(segment) > 1 and np.std(segment) > 1e-6:
            sr = np.mean(segment) / np.std(segment)
        else:
            sr = np.nan
        row[period_label] = round(sr, 4)
    table4_rows.append(row)

table4 = pd.DataFrame(table4_rows).set_index('Strategy')
print("TABLE 4 | Sharpe Ratio (4-month window)")
display(table4)
\end{lstlisting}
\end{tcolorbox}

\vspace{6pt}
\begin{table}[H]
  \centering
  \caption{TABLE 4 --- Sharpe Ratio (jendela 4 bulan)}
  \small
  \begin{tabular}{lrrrrrr}
    \toprule
    \textbf{Strategy} & \textbf{Jan-18} & \textbf{May-18} & \textbf{Sep-18} & \textbf{Jan-19} & \textbf{May-19} & \textbf{Sep-19} \\
    \midrule
    \midrule
    EW                        & -0.2532 & -0.0425 & -0.0986 & -0.1270 &  0.2016 & -0.1298 \\
    CM                        & -0.2508 & -0.0984 & -0.0577 & -0.1500 &  0.1904 & -0.0598 \\
    GM                        & -0.2355 & -0.1032 & -0.0537 & -0.1873 &  0.1775 & -0.0526 \\
    NW (gamma=1.0)            & -0.3849 & -0.0450 & -0.0632 & -0.0152 &  0.1567 & -0.0498 \\
    AGGP (Optimized)          & -0.3303 & -0.0552 & -0.0794 & -0.1797 &  0.1635 & -0.0902 \\
    \bottomrule
  \end{tabular}
\end{table}

\textbf{Analisis Tabel 4:} \textit{Sharpe Ratio} yang negatif di hampir semua periode (kecuali Mei 2019) mencerminkan kondisi pasar yang didominasi oleh tren turun. Meskipun demikian, pada periode pemulihan (Mei 2019), hampir seluruh strategi memberikan \textit{risk-adjusted return} yang positif. Strategi berbasis graf cenderung memiliki Sharpe Ratio yang lebih kompetitif dalam periode transisi dibandingkan Markowitz klasik yang sering kali terlalu sensitif terhadap estimasi parameter yang tidak stabil.

% ============================================================================
\section{Sel 14 --- Analisis Tail-Risk (Table 5 --- Rachev Ratio)}
% ============================================================================

Rachev Ratio (RR) memberikan wawasan tentang asimetri distribusi imbal hasil
pada ekor (\emph{tails}):
\[
  \text{RR} = \frac{\text{CVaR}_{\text{upper}}(10\%)}{\text{CVaR}_{\text{lower}}(10\%)}
\]

\cellIn{15}
\begin{tcolorbox}[codebox]
\begin{lstlisting}[language=Python, numbers=none]
table5_rows = []
for strat_name, res in results.items():
    df_res = res['results_df'].set_index('date')
    row    = {'Strategy': strat_name}
    for period_label, start, end in period_ranges:
        segment = df_res.loc[start:end, 'return'].values
        if len(segment) > 5:
            rr = calculate_rachev_ratio(segment, alpha=0.10)
        else:
            rr = np.nan
        row[period_label] = round(rr, 4)
    table5_rows.append(row)

table5 = pd.DataFrame(table5_rows).set_index('Strategy')
print("TABLE 5 | Rachev Ratio (alpha=10%, 4-month window)")
display(table5)
\end{lstlisting}
\end{tcolorbox}

\vspace{6pt}
\begin{table}[H]
  \centering
  \caption{TABLE 5 --- Rachev Ratio ($\alpha=10\%$, jendela 4 bulan)}
  \small
  \begin{tabular}{lrrrrrr}
    \toprule
    \textbf{Strategy} & \textbf{Jan-18} & \textbf{May-18} & \textbf{Sep-18} & \textbf{Jan-19} & \textbf{May-19} & \textbf{Sep-19} \\
    \midrule
    \midrule
    EW                        & 0.3926 & 0.8889 & 0.8062 & 0.7200 & 1.7430 & 0.5381 \\
    CM                        & 0.3969 & 0.7612 & 0.7865 & 0.6506 & 2.0115 & 0.7510 \\
    GM                        & 0.4126 & 0.7515 & 0.7760 & 0.5600 & 1.6893 & 0.7353 \\
    NW (gamma=1.0)            & 0.3488 & 1.0222 & 0.8170 & 1.0171 & 2.1687 & 0.8614 \\
    AGGP (Optimized)          & 0.5188 & 1.0019 & 0.7941 & 0.5625 & 1.8104 & 0.6928 \\
    \bottomrule
  \end{tabular}
\end{table}

\textbf{Analisis Tabel 5:} \textit{Rachev Ratio} memberikan gambaran asimetri risiko. Pada Januari 2019, strategi \textit{Network Markowitz} dan diversifikasi graf menunjukkan RR $> 1$, yang berarti potensi imbal hasil di ekor kanan (\textit{upside}) mulai mengimbangi risiko di ekor kiri (\textit{downside}). \textbf{AGGP} menunjukkan performa yang sangat kuat pada awal periode (Jan-2018) dengan RR 0.51, mengungguli NW dan strategi diversifikasi lainnya dalam hal ketahanan terhadap \textit{tail-risk}.

% ============================================================================
\section{Sel 15 --- Analisis Fase Pasar (Table 6)}
% ============================================================================

Pemetaan performa ke tiga fase utama pasar:

\begin{itemize}
  \item \textbf{Bearish}: Januari 2018 -- Maret 2019
  \item \textbf{Recovery}: April 2019 -- Juni 2019
  \item \textbf{Stable/Sideways}: Juli 2019 -- Oktober 2019
\end{itemize}

\cellIn{16}
\begin{tcolorbox}[codebox]
\begin{lstlisting}[language=Python, numbers=none]
fase_pasar = {
    'Bearish':  ('2018-01-01', '2019-03-31'),
    'Recovery': ('2019-04-01', '2019-06-30'),
    'Stable':   ('2019-07-01', '2019-10-17')
}

sharpe_rows = []
rachev_rows = []
for strat_name, res in results.items():
    df_res = res['results_df'].set_index('date')
    sr_row = {'Strategy': strat_name}
    rr_row = {'Strategy': strat_name}
    for fase, (start, end) in fase_pasar.items():
        segment = df_res.loc[start:end, 'return'].values
        # Sharpe Ratio
        if len(segment) > 1 and np.std(segment) > 1e-6:
            sr = np.mean(segment) / np.std(segment)
        else:
            sr = np.nan
        sr_row[fase] = round(sr, 4)
        # Rachev Ratio
        if len(segment) > 5:
            rr = calculate_rachev_ratio(segment, alpha=0.10)
        else:
            rr = np.nan
        rr_row[fase] = round(rr, 4)
    sharpe_rows.append(sr_row)
    rachev_rows.append(rr_row)

table6_sharpe = pd.DataFrame(sharpe_rows).set_index('Strategy')
table6_rachev = pd.DataFrame(rachev_rows).set_index('Strategy')

print("TABLE 6 | Market Phase Performance")
print("\nPanel A: Sharpe Ratio")
display(table6_sharpe)
print("\nPanel B: Rachev Ratio (alpha=10%)")
display(table6_rachev)
\end{lstlisting}
\end{tcolorbox}

\vspace{6pt}
\begin{table}[H]
  \centering
  \caption{TABLE 6 --- Panel A: Sharpe Ratio per Fase Pasar}
  \begin{tabular}{lrrr}
    \toprule
    \textbf{Strategy} & \textbf{Bearish} & \textbf{Recovery} & \textbf{Stable} \\
    \midrule
    \midrule
    EW                        & -0.0726 &  0.1286 & -0.1301 \\
    CM                        & -0.0923 &  0.1844 & -0.1022 \\
    GM                        & -0.0991 &  0.1732 & -0.0950 \\
    NW (gamma=1.0)            & -0.0685 &  0.1359 & -0.0628 \\
    AGGP (Optimized)          & -0.1081 &  0.1690 & -0.1088 \\
    \bottomrule
  \end{tabular}
\end{table}

\begin{table}[H]
  \centering
  \caption{TABLE 6 --- Panel B: Rachev Ratio per Fase Pasar ($\alpha=10\%$)}
  \begin{tabular}{lrrr}
    \toprule
    \textbf{Strategy} & \textbf{Bearish} & \textbf{Recovery} & \textbf{Stable} \\
    \midrule
    \midrule
    EW                        & 0.7933 & 1.2168 & 0.5712 \\
    CM                        & 0.6463 & 1.2463 & 0.7523 \\
    GM                        & 0.6275 & 1.1627 & 0.7477 \\
    NW (gamma=1.0)            & 0.7244 & 1.4069 & 0.8656 \\
    \textbf{AGGP (Optimized)} & 0.7144 & \textbf{1.6631} & 0.6830 \\
    \bottomrule
  \end{tabular}
\end{table}

\textbf{Analisis Tabel 6:} Segmentasi per fase pasar mengonfirmasi bahwa \textbf{AGGP} adalah model yang paling "tahan banting". Pada fase \textit{Recovery}, AGGP mencatatkan \textit{Rachev Ratio} tertinggi sebesar \textbf{1.6631}, yang menunjukkan kemampuan adaptasi yang luar biasa saat pasar berbalik arah. Sementara itu, pada fase \textit{Bearish}, strategi NW dengan penalti tinggi ($\gamma=1.0$) memberikan kestabilan yang lebih baik dibandingkan strategi Markowitz standar, memperkuat argumen bahwa struktur jaringan memberikan informasi penting yang tidak tertangkap oleh korelasi historis biasa.

% ============================================================================
\section{Ringkasan dan Temuan Utama}
% ============================================================================

Berdasarkan seluruh eksperimen yang telah dijalankan, beberapa temuan utama
dapat dirangkum sebagai berikut:

\begin{enumerate}[leftmargin=2cm]
  \item \textbf{AGGP (Optimized)} secara konsisten unggul dalam hal kerugian
        kumulatif terendah di semua periode, dengan hanya $-36.22\%$ pada
        September 2019 dibandingkan $-88.82\%$ untuk EW.

  \item \textbf{Network Markowitz dengan $\gamma$ tinggi} ($\gamma \geq 0.7$)
        menunjukkan keunggulan dibandingkan NW tanpa penalti, mengonfirmasi
        bahwa penalti sentralitas efektif menekan eksposur ke aset yang sangat
        terhubung dalam jaringan.

  \item \textbf{Dynamic Graph Diversification} justru menghasilkan performa
        paling buruk ($-91.22\%$), menunjukkan bahwa optimasi threshold via
        cross-validation sangat rentan terhadap \emph{overfitting} pada data
        kripto yang sangat volatil.

  \item Pada fase \textbf{Recovery}, hampir semua strategi memiliki Rachev
        Ratio $> 1$, yang mengindikasikan bahwa \emph{upside tail} lebih besar
        dari \emph{downside tail} --- sebuah kondisi menguntungkan.

  \item \textbf{AGGP} mencatat Rachev Ratio tertinggi pada fase Recovery
        (1.6631), mengonfirmasi kemampuannya sebagai \emph{all-terrain model}
        yang adaptif terhadap berbagai kondisi pasar.
\end{enumerate}

\textbf{Simpulan Validitas Akademik:}
Studi ini menunjukkan bahwa model AGGP memenuhi kaidah \textit{fairness} akademik melalui pengujian \textit{Out-of-Sample} yang ketat dan inklusi biaya transaksi yang realistis. Superioritas model ini bukan merupakan hasil dari \textit{data snooping}, melainkan berasal dari integrasi informasi topologi jaringan yang mampu menangkap dinamika risiko sistemik yang sering terlewatkan oleh metrik keuangan tradisional. Hal ini memvalidasi penggunaan teori graf sebagai instrumen pelengkap yang krusial bagi sistem konsultasi keuangan otomatis (*Robo-Advisory*) di masa depan.

% ============================================================================
\section*{Referensi}
\addcontentsline{toc}{section}{Referensi}
% ============================================================================

\begin{itemize}
  \item Giudici, P., Sariev, A., \& Toscani, G. (2020).
        \textit{Network Models to Improve Automated Cryptocurrency Portfolio
        Management}. Risks, 8(3), 96.
        \url{https://doi.org/10.3390/risks8030096}

  \item Marchenko, V.~A., \& Pastur, L.~A. (1967).
        Distribution of eigenvalues for some sets of random matrices.
        \textit{Mathematics of the USSR-Sbornik}, 1(4), 457--483.

  \item Mantegna, R.~N. (1999).
        Hierarchical structure in financial markets.
        \textit{The European Physical Journal B}, 11(1), 193--197.
\end{itemize}

\end{document}
